\subsection{Further strong-component reductions for the baseline}

\begin{lemma}[Proved: a binary idempotent representable]
\label{CP2:A-baseline-binary-representable}
Let $S$ be a finite category, let $e\in S(a,a)$ be idempotent,
and suppose $1\leq |S(a,i)e|\leq2$ for every object $i$.
Every positive-degree rational cohomology product of $BS$
is zero.  There is a dual assertion for $eS(i,a)$.
\end{lemma}
\begin{proof}
Let $F(i)=S(a,i)e$, as a covariant finite act.  The category
of elements of the full representable $S(a,-)$ has initial
object $(a,1_a)$ and is contractible.  The maps $u\mapsto ue$
define a natural idempotent endomorphism of this act and
give a retraction of its category of elements onto the
category of elements of $F$.  Its classifying space is
therefore a retract of a contractible space and is itself
contractible: compose a contraction with the inclusion and
retraction.  The act cone of $F$ is the homotopy cofiber of
this contractible space mapping to $BS$, and the inclusion
of $BS$ in the cofiber is a homotopy equivalence.  The
left-act version of Lemma~\ref{CP2:A-high-affine} applies
to the one- or two-element stalks and annihilates all
positive rational products.  Taking opposites proves the
dual statement.
\end{proof}

\begin{lemma}[Proved: an objectwise binary projective cover]
\label{CP2:A-baseline-binary-projective-cover}
Let $S$ be a finite category.  Suppose the constant left
$\mathbb QS$-module $\mathbf1$ is the augmentation quotient
of a projective module $P$ which at each object has a basis
of one or two elements, and whose arrow maps send basis
elements to basis elements.  Every finite right $S$-act
with all stalks nonempty has an act cone with zero rational
degree-two products.  An opposite endpoint whose trivial
right rational module is projective may replace the trivial
sink.  There are left-right duals.
\end{lemma}
\begin{proof}
Put $R=\mathbb QS$ and $V=\ker(P\to\mathbf1)$.
Each $V_i$ has dimension zero or one.  Use the difference
of the two basis vectors when it has dimension one.
For an arrow $f:i\to j$ the induced difference map is
the scalar $\chi_f\in\{0,1,-1\}$, with zero whenever
one difference space is zero.  These typed maps compose
as module maps.

For the right act $F$, set
$L=\ker(\mathbb QF\to\mathbf1_{\rm right})$, using
objectwise augmentation.  All stalks are nonempty, so
this sequence is exact.  Tensoring a projective left
resolution with it, as in the relative calculation in
Lemma~\ref{CP2:A-baseline-eight-act}, identifies reduced
act-cone chains with $(L\otimes_R^{\mathbf L}\mathbf1)[1]$.
Finite-dimensional rational duality and the projectivity
of $P$ give a natural surjection
\[
 \operatorname{Hom}_R(V,DL)\longrightarrow
 \operatorname{Ext}^1_R(\mathbf1,DL)
  =H^2(B\mathcal C(S,F);\mathbb Q).
\]
Here $D$ denotes rational vector-space duality, object by
object, and the first map is the connecting homomorphism
of the projective-cover sequence.

Lift a second class through this map.  In the chosen
difference coordinates it is a family of functionals
$\ell_i:L_i\to\mathbb Q$, zero when $V_i=0$, satisfying
$\ell_i L(f)=\chi_f\ell_j$.  Choose $s_i\in F(i)$ and
put $q_i(s)=\ell_i([s]-[s_i])$.  Subtracting two such
values gives
\[
 q_i(F(f)s)=\chi_f q_j(s)+\beta_f,
 \qquad \beta_f=q_i(F(f)s_j).
\]
Thus the chosen class is encoded by a finite affine
diagram of rational values.

Take one rational grid spacing $\delta>0$ containing
all these values and translations.  Let $I_i$ be the
convex hull of $q_i(F(i))$ and let
$T_i=I_i\cap\delta\mathbb Z$.  The displayed affine
maps preserve the intervals and their grids.  They
compose exactly: on a nondegenerate source interval
their values at the two endpoints determine them, and
on a degenerate interval there is only one value to
check.  Therefore $T$ is an actual finite right act,
and $q:F\to T$ is an actual act map.

On $\ker(\mathbb QT_i\to\mathbb Q)$, the coordinate
functional $\sum a_t[t]\mapsto\sum a_t t$ satisfies
the same typed equation: translations disappear because
$\sum a_t=0$.  It is zero when $V_i=0$, and its pullback
is $\ell_i$.  Naturality of the connecting homomorphism
therefore lifts the original class from a second class
on the act cone of $T$.

For clarity, the common-grid extension of the interval
argument is strict.  Each $I_i/T_i$ is a wedge of circles
indexed by consecutive grid intervals of length $\delta$.
An affine map of slope $1$ or $-1$ takes one grid interval
onto another of the same length, preserving or reversing
its circle parameter; slope zero gives a pointed constant.
The degree-two circle maps commute with every such map.
They induce a self-map of the pointed category bar which
is multiplication by two on its reduced cellular chain
complex, in every total degree.  It is therefore
multiplication by two on every positive cohomology group.
Multiplicativity gives $2uv=4uv$ for positive classes,
so all those products vanish rationally.  In particular
the square of the lifted class, and hence of the original
class, is zero.  Rational polarization proves the mixed
degree-two assertion.

For a projective opposite endpoint, the exact orbit
comparison of Lemma~\ref{CP2:A-high-orbits} gives an
actual functor inducing a rational cohomology-ring
isomorphism to an act cone.  Its stalks remain nonempty,
so the preceding argument applies.  Opposition gives
the remaining statements.
\end{proof}

\begin{corollary}[Proved: two-row and two-orbit instances]
\label{CP2:A-baseline-two-row-instances}
The preceding lemma applies when $S$ is strongly connected
with trivial minimal group and at most two minimal rows
at every object.  It also applies whenever there are
$a,e,G$, with $e^2=e$ and $G$ a finite subgroup of the
units of $eS(a,a)e$, such that every orbit set
$S(a,i)e/G$ has one or two elements.
\end{corollary}
\begin{proof}
In the first case take $P=(\mathbb QS)e$ at a minimal
idempotent.  Its bases are the minimal rows at each object.
In the second case take
$P=(\mathbb QS)p$, where $p=|G|^{-1}\sum_{g\in G}g$.
It is projective because $p$ is idempotent.  The normalized
averages of the $G$-orbits give the required basis.  An
equivariant map sends a normalized orbit average to the
normalized average of its image orbit: all fibers have
equal size.  Thus every arrow acts by a basis map, and
objectwise augmentation onto the constant module is
surjective.  These are precisely the hypotheses needed.
\end{proof}

\begin{proposition}[Proved: wholly strong two-object categories through seventeen]
\label{CP2:A-baseline-strong-two}
A strongly connected category on two objects with at most
seventeen arrows cannot realize $\mathbb{CP}^2$.
\end{proposition}
\begin{proof}
If one object is a retract of the other, deletion saves
at least three arrows and leaves a monoid of order at most
fourteen.  Theorem~\ref{A:monoid-minimum23-classification}
excludes it.  Suppose no object retract is present.
Both identities lie outside the categorical minimal ideal.
Write $r=r_0+r_1$ and $c=c_0+c_1$ for its total minimal
row and column counts and $G$ for its maximal group.
If $r=2$ or $c=2$, all local row or column counts are one
and the category has finite-group homotopy type by
Lemma~\ref{CP2:A-high-Rees}; its positive rational
cohomology is zero.

Otherwise $r,c\geq3$ and $|G|rc+2\leq17$.
Consequently $G=1$ and $\min(r,c)=3$, since
$2\cdot3\cdot3+2=20$ and $4\cdot4+2=18$.
On the side with total three the two positive local
counts are one and two.  The corresponding idempotent
representable has one or two states at each object.
Lemma~\ref{CP2:A-baseline-binary-representable}, or its
dual, makes every positive rational product zero.
The nonzero projective-plane square is impossible.
\end{proof}

\begin{proposition}[Proved: cyclic-support three-object reduction through sixteen]
\label{CP2:A-baseline-strong-three}
A three-object category with cyclic support and at most
sixteen arrows either fails to realize $\mathbb{CP}^2$,
or admits an object-retract deletion to a two-object
acyclic-support category with at most thirteen arrows.
In particular, the cyclic-support part of the
three-object baseline requires only that explicit
two-object baseline dependency.
\end{proposition}
\begin{proof}
If the category is wholly strong, a retract deletion
leaves a wholly strong two-object category and the
preceding proposition applies.  Without an object
retract, its three identities are outside the minimal
ideal.  Unless it has finite-group type, the total
row and column counts are at least four, requiring
at least $4\cdot4+3=19$ arrows.  Thus the wholly
strong case is excluded.

The remaining connected case has a proper strong pair
$S$ and one object with corner $N$.  A retract inside
the pair has exactly the conclusion asserted: deletion
saves at least three arrows and the remaining support
is acyclic.  Assume that no object retract remains.
By opposition take $S$ as the source component.  There
are nonempty outgoing stalks at both its objects.

If $H^2(BS;\mathbb Q)$ and $H^2(BN;\mathbb Q)$ both
vanish, use the two-block lemma.  If the first is
nonzero, its total minimal rows and columns are at
least three.  Since $|S|\leq16-1-2=13$, the minimal
group is trivial, both totals are exactly three,
and $|S|\geq11$.  Every local row count is at most
two.  The other corner has order at most three.
If its order is at most two, its trivial right module
is projective and the two-row instance of the preceding
lemma kills all second products on the collage.  If
its order is three, the entire budget forces
$|S|=11$ and exactly one outgoing arrow at each object.
The collage is the full categorical join of $S$ and
$N$, whose positive cohomology products vanish.

It remains that $H^2(BS;\mathbb Q)=0$ and
$H^2(BN;\mathbb Q)\ne0$.  Then $5\leq|N|\leq8$
and $|S|\leq9$.  The two identities of $S$ are
outside its ideal.  Its minimal group is trivial,
since a group of order at least two would already
cost $2\cdot2\cdot2+2=10$ internal arrows.  One
total row or column count must be two, because
$3\cdot3+2=11>9$.

If the row total is two, the one-row source-component
collapse of Lemma~\ref{CP2:A-high-one-sided} gives an
ordinary left-act cone over $N$.  This is an actual
homotopy equivalence, so its rational homology is
bounded under the putative projective-plane hypothesis.
Lemma~\ref{CP2:A-baseline-eight-act} excludes it.
Otherwise the column total is two and the row total
is three, so $|S|\geq8$.  If its two outgoing stalks
are singleton, the join argument applies.  If not,
the only budget is $|S|=8$, $|N|=5$, and three
outgoing arrows.  The pair consists of its six-arrow
minimal band and its two identities; its local row
counts are two and one, and both local column counts
are one.  At the one-row object every incoming
internal arrow is fixed by its local minimal
idempotent, and its only unit is its identity.
There are no external entrances because $S$ is a
source.  The one-sided object deletion therefore
applies.  The retained collage has at most two
cross arrows, since either retained outgoing stalk
had size one or two.  It is excluded by the
singleton join or the universal two-cross lemma.
This exhausts the proper-pair cases without using
an inherited composition census.
\end{proof}
